\documentclass[11pt]{article}
\usepackage[margin=1in]{geometry}
\usepackage{enumitem}
% =============================================================================
% Preambles/header.tex — shared LaTeX/Beamer preamble for this project
%
% Use from a Beamer lecture by prepending:
%     \input{header}
% and compiling with TEXINPUTS=../Preambles:$TEXINPUTS (the /compile-latex
% skill does this automatically).
%
% PALETTE CONTRACT. Color names defined here MUST match the names in
% Quarto/theme-template.scss so Beamer and Quarto renderings use the same
% palette. The scripts/check-palette-sync.sh script enforces this.
%
% To customize: edit the HEX values below AND the matching $variable in
% Quarto/theme-template.scss, then run scripts/check-palette-sync.sh.
% =============================================================================

% -----------------------------------------------------------------------------
% Engine detection + encoding
%
% Our /compile-latex skill uses XeLaTeX, where inputenc is unnecessary and
% can produce encoding warnings. Only load inputenc under pdfTeX.
% -----------------------------------------------------------------------------
\usepackage{iftex}
\ifPDFTeX
  \usepackage[utf8]{inputenc}
\fi

\usepackage{amsmath, amssymb, amsthm}
\usepackage{booktabs}
\usepackage{graphicx}

% -----------------------------------------------------------------------------
% PALETTE — mirrors Quarto/theme-template.scss variable names.
% Load xcolor and define colors BEFORE anything that references them
% (notably hyperref, hypersetup, and the Beamer color assignments).
% -----------------------------------------------------------------------------
\usepackage{xcolor}

\definecolor{primary-blue}{HTML}{012169}      % headings, accents
\definecolor{primary-gold}{HTML}{B9975B}      % emphasis, borders
\definecolor{highlight-yellow}{HTML}{F2A900}  % markers, alerts
\definecolor{light-bg}{HTML}{E8EDF5}          % title / section backgrounds
\definecolor{jet}{HTML}{1A1A1A}               % body text

% Semantic colors (used in diagrams and annotations)
\definecolor{positive}{HTML}{15803D}          % good / observed / identified
\definecolor{negative}{HTML}{B91C1C}          % bad / problematic / bias
\definecolor{neutral}{HTML}{525252}           % reference / context

% Highlight accents (match .hi-* classes in the SCSS)
\definecolor{hi-slate}{HTML}{314F4F}
\definecolor{hi-green}{HTML}{2E8B57}
\definecolor{hi-red}{HTML}{C41E3A}

% Semantic aliases so skill templates and snippets can write
% \textcolor{accent}{...} without hard-coding "primary-blue".
\colorlet{accent}{primary-blue}
\colorlet{accent2}{primary-gold}

% -----------------------------------------------------------------------------
% Hyperref — loaded AFTER xcolor + palette so link colors resolve.
% -----------------------------------------------------------------------------
\usepackage{hyperref}
\hypersetup{colorlinks=true, linkcolor=primary-blue, urlcolor=primary-blue,
            citecolor=primary-blue, pdfborder={0 0 0}}

% -----------------------------------------------------------------------------
% Course code — set once per deck (after \input{header}, before \begin{document})
% via \coursecode{CS401}. Shown in the footer on every slide so a deck's course
% is identifiable without opening the title page. Empty by default.
% -----------------------------------------------------------------------------
\makeatletter
\newcommand{\coursecode}[1]{\gdef\@coursecodetext{#1}}
\gdef\@coursecodetext{}
\makeatother

% -----------------------------------------------------------------------------
% Beamer theme assignments (only apply under Beamer).
%
% We detect Beamer via \@ifclassloaded{beamer}, which is the documented,
% non-internal way. This must be wrapped in \makeatletter/\makeatother
% because \@ifclassloaded contains an @-catcode character.
% -----------------------------------------------------------------------------
\makeatletter
\@ifclassloaded{beamer}{%
  \usefonttheme{professionalfonts}
  \setbeamercolor{structure}{fg=primary-blue}
  \setbeamercolor{frametitle}{fg=primary-blue}
  \setbeamercolor{title}{fg=primary-blue}
  \setbeamercolor{subtitle}{fg=primary-gold}
  \setbeamercolor{author}{fg=primary-blue}
  \setbeamercolor{institute}{fg=neutral}
  \setbeamercolor{date}{fg=primary-gold}
  \setbeamercolor{itemize item}{fg=primary-blue}
  \setbeamercolor{itemize subitem}{fg=primary-gold}
  \setbeamercolor{alerted text}{fg=primary-gold}
  \setbeamercolor{block title}{fg=white, bg=primary-blue}
  \setbeamercolor{block body}{bg=light-bg}
  % Semantic block variants: block=definition/concept (blue), exampleblock=
  % worked example (green/positive), alertblock=key takeaway or caution (gold).
  % Keep to <=2 colored boxes per slide across ALL three variants (INV-7).
  \setbeamercolor{block title example}{fg=white, bg=positive}
  \setbeamercolor{block body example}{bg=light-bg}
  \setbeamercolor{block title alerted}{fg=white, bg=primary-gold}
  \setbeamercolor{block body alerted}{bg=light-bg}

  % Minimal footer: course code (left) + page X / N (right), no navigation clutter
  \setbeamertemplate{navigation symbols}{}
  \setbeamertemplate{footline}{%
    \color{neutral}\footnotesize\hspace{0.7em}\@coursecodetext\hfill
    \insertframenumber/\inserttotalframenumber\hspace{0.7em}\vspace{0.3em}}
}{}
\makeatother

% -----------------------------------------------------------------------------
% TikZ setup — libraries the snippet gallery uses
% -----------------------------------------------------------------------------
\usepackage{tikz}
\usetikzlibrary{arrows.meta, positioning, calc, decorations.pathreplacing,
                fit, shapes.geometric, backgrounds}

% Shared TikZ styles — snippets and hand-written diagrams can pick these up
% via `\begin{tikzpicture}[every node/.style={...}]` or by naming specific
% styles (e.g., `\node[dag-node] (X) at (0,0) {X};`).
\tikzset{
  % Node styles for causal DAGs and similar
  dag-node/.style = {
    draw=primary-blue, fill=light-bg, rounded corners=2pt,
    minimum width=1.2cm, minimum height=0.9cm, align=center, font=\small
  },
  decision-node/.style = {
    draw=primary-gold, fill=white, diamond, aspect=1.8,
    minimum width=1.8cm, minimum height=1.2cm, align=center, font=\small,
    inner sep=1pt
  },
  % Edge styles — observed vs counterfactual
  observed-edge/.style = {-{Latex[length=2.5mm]}, thick, draw=jet},
  counterfactual-edge/.style = {-{Latex[length=2.5mm]}, thick, draw=neutral, dashed},
  confound-edge/.style = {-{Latex[length=2.5mm]}, thick, draw=negative, dashed},
  % Data-point styles for plots
  observed-dot/.style = {fill=primary-blue, draw=primary-blue, circle, inner sep=1.2pt},
  counterfactual-dot/.style = {fill=white, draw=neutral, circle, inner sep=1.2pt}
}

% -----------------------------------------------------------------------------
% Convenience macros
% -----------------------------------------------------------------------------
\newcommand{\muted}[1]{\textcolor{neutral}{#1}}
\newcommand{\key}[1]{\textcolor{primary-gold}{\textbf{#1}}}
\newcommand{\good}[1]{\textcolor{positive}{#1}}
\newcommand{\bad}[1]{\textcolor{negative}{#1}}

% Standout frame macro: full-bleed dark background for section transitions.
% Beamer-only (uses \begin{frame}/\setbeamercolor/\usebeamerfont) -- do not
% call from an article-class file (e.g. Notes/<CODE>/*-notes.tex). \muted,
% \key, \good, \bad, and \coursecode above are plain xcolor/gdef and are
% safe to use from either class; verified when Notes/article-class support
% was added.
% Expand: \transitionslide{Your transition title}
\newcommand{\transitionslide}[1]{%
  \begingroup
  \setbeamercolor{background canvas}{bg=primary-blue}
  \begin{frame}[plain]
    \centering
    \vfill
    {\usebeamerfont{title}\color{white}\Large #1\par}
    \vfill
  \end{frame}
  \endgroup
}

% Section divider frame (Beamer-only), an alternative to \transitionslide with a
% darker two-line style: near-black (jet) background, a small white label line
% above a larger gold title, and a thin gold rule along the bottom edge.
% Expand: \sectiondivider{Section 2}{Pointers}
\newcommand{\sectiondivider}[2]{%
  \begingroup
  \setbeamercolor{background canvas}{bg=jet}
  \begin{frame}[plain]
    \vspace*{3.0cm}
    {\color{white}\ttfamily\large #1\par}
    \vspace{0.5cm}
    {\color{primary-gold}\LARGE #2\par}
    \begin{tikzpicture}[remember picture, overlay]
      \draw[primary-gold, line width=2.5pt]
        ($(current page.south west)+(0,0.1)$) -- ($(current page.south east)+(0,0.1)$);
    \end{tikzpicture}
  \end{frame}
  \endgroup
}

% -----------------------------------------------------------------------------
% Channel watermark — "Concepts That Click" (YouTube).
%
% Article-class files (CompetitiveExam Q/A sets, Notes, Assignments) get a
% light diagonal draft watermark on every page plus a centered footer naming
% the channel. Beamer decks get a subtle TikZ background watermark instead
% (the footline already carries the course code). Centralized here so every
% MCQ PDF is branded without per-file edits.
% -----------------------------------------------------------------------------
\makeatletter
\@ifclassloaded{article}{%
  \usepackage{draftwatermark}
  \SetWatermarkText{Concepts That Click}
  \SetWatermarkScale{1}
  \SetWatermarkFontSize{2cm}
  \SetWatermarkLightness{0.86}
  \SetWatermarkAngle{45}
  \usepackage{fancyhdr}
  \pagestyle{fancy}
  \fancyhf{}
  \fancyfoot[C]{\footnotesize\textcolor{neutral}{Concepts That Click~|~YouTube: \href{https://www.youtube.com/@concepts.thatclick}{@concepts.thatclick}}}
  \renewcommand{\headrulewidth}{0pt}
  \renewcommand{\footrulewidth}{0pt}
  \fancypagestyle{plain}{%
    \fancyhf{}%
    \fancyfoot[C]{\footnotesize\textcolor{neutral}{Concepts That Click~|~YouTube: \href{https://www.youtube.com/@concepts.thatclick}{@concepts.thatclick}}}%
    \renewcommand{\headrulewidth}{0pt}%
    \renewcommand{\footrulewidth}{0pt}%
  }
}{}
\@ifclassloaded{beamer}{%
  \setbeamertemplate{background}{%
    \begin{tikzpicture}[remember picture, overlay]
      \node[rotate=45, opacity=0.055, align=center]
        at (current page.center)
        {\Huge\textcolor{neutral}{Concepts That Click}};
    \end{tikzpicture}%
  }
}{}
\makeatother 
\coursecode{JKSSB}
\renewcommand{\thesection}{48.\arabic{section}}
\newtheorem{example}{Example}[section]
\newenvironment{definitionbox}[1]{%
  \par\noindent\textbf{Definition (#1).}\ \itshape
}{\par\vspace{0.3em}}
\title{MCQ Masterclass --- Detailed Lecture Notes}
\author{JKSSB Computer Awareness}
\date{\today}

\begin{document}
\maketitle

\section{How to use this capstone}

This lesson is a rapid-revision drill over the whole Computer Awareness
course. It does not introduce new topics. It collects, unit by unit, the
must-remember ladder (the facts worth one mark each) and the trap catalogue
(the wrong beliefs the distractors rely on). Read each section as a pair:
first the ladder sentence, then the trap that inverts it. Where a trap has a
course registry ID (TRAP-01 and so on), the ID is given so the item can be
traced back to the evidence.

\begin{center}
\begin{tabular}{p{0.24\textwidth}p{0.66\textwidth}}
\toprule
\textbf{Term} & \textbf{Meaning in this lesson} \\
\midrule
Must-remember ladder & The one-line fact per unit that directly answers a recall item. \\
Trap catalogue & The near-neighbour wrong belief the examiner plants as a distractor. \\
Near-neighbour contrast & The discipline of stating what a term is, what it is not, and its nearest confusion. \\
\bottomrule
\end{tabular}
\end{center}

\section{Memory and storage ladder}

The speed order, fastest first, is registers, cache (L1, L2, L3), main memory
(RAM), secondary storage (SSD/HDD), then tertiary or offline storage. RAM
(Random Access Memory) is volatile primary memory: its contents are lost when
power goes off. ROM (Read-Only Memory) is non-volatile and read-only in
normal operation; it holds firmware such as the BIOS startup instructions.
Cache is faster than main memory and holds recently and frequently used data;
it keeps copies of data that live in RAM and is therefore not a backup. The
CPU works directly with main memory and does not access secondary storage
directly.

\begin{definitionbox}{Volatile vs non-volatile}
Volatile memory loses its contents without power (RAM, cache, registers).
Non-volatile memory retains its contents without power (ROM, SSD, HDD,
magnetic tape).
\end{definitionbox}

The traps: RAM is non-volatile (TRAP-01), ROM allows normal read and write
(TRAP-02), and the CPU accesses secondary storage directly (TRAP-03) are all
false. Cache being slower than RAM, or cache being a backup, inverts the same
ladder.

\section{Hardware, software and device ladder}

A device driver is software, and the operating system is system software
(TRAP-10). Open-source means the licence permits viewing, changing, and
sharing the source; it is not the same as freeware (free of cost) or shareware
(trial) (TRAP-23). The firmware that initialises the hardware and loads the
operating system is the BIOS, Basic Input/Output System (TRAP-24). The control
unit controls the operation sequence and data flow. Input devices in the
exam's drill set are the light pen, punch-card reader, digitizer tablet, and
scanner; output devices are the plotter, drum printer, and sound card
(TRAP-12). MICR reads magnetic-ink characters, OCR reads printed text, and OMR
reads marks (TRAP-25). Large engineering drawings call for a plotter; ICs
arrived in the third generation and microprocessors in the fourth (TRAP-26).
The smallest unit of data is the bit; 8 bits make one byte (TRAP-27).

\begin{example}
An item lists plotter, scanner, sound card, and digitizer tablet and asks which
is input. Trace the drill: plotter is output, sound card is output, scanner is
input, digitizer tablet is input. Any option calling the plotter an input
device is the planted distractor.
\end{example}

\section{Operating system ladder}

The operating system is the layer that provides a user-friendly interface
between the user and the hardware. While it runs, it is resident in primary
storage (main memory). Multitasking means running more than one program at a
time. Android is the open-source mobile OS of the syllabus. These four
sentences answer four distinct direct-recall items, so memorise them verbatim.

\section{MS Word, Excel, Access and PowerPoint ladder}

In MS Word, F7 runs spell check and Mail Merge combines letters with an
address database (TRAP-29). A Section Break, not a Page Break, enables
different headers and footers in different parts of a document (TRAP-13). In
MS Excel, every formula must start with \texttt{=}; \texttt{\$A\$1} is an
absolute reference that does not shift when copied; COUNTIFS counts rows
meeting several criteria, and SUMPRODUCT with multiplication does equivalent
multi-criterion counting (TRAP-14, TRAP-30). In MS Access, the Select query
retrieves data on a condition, while Update, Append, and Delete modify data
(TRAP-31). In PowerPoint, F5 starts the slide show from the beginning; Shift
plus F5 starts from the current slide (TRAP-32). The string \texttt{.png} is an
extension; PNG is the format --- never conflate the two (TRAP-34).

\begin{example}
An item asks which break allows different headers in the two halves of a
document. The Page Break option is the trap: a Page Break only pushes content
to a new page. Only a Section Break restarts headers and footers, so it is the
correct choice.
\end{example}

\section{Internet, web and email ladder}

The Internet is the global network of networks; the WWW (World Wide Web) is
the system of linked pages served over it. An ISP --- Internet Service
Provider --- provides the connection. A URL (Uniform Resource Locator) is the
address of a resource, and DNS (Domain Name System) resolves names to IP
addresses. IPv4 is 32-bit; IPv6 is 128-bit with a fixed 40-byte header
(TRAP-18). In an email address, \texttt{@} separates the user name from the
domain name (TRAP-28). BCC hides recipients from each other while the sender
always knows and To/Cc stay visible (TRAP-16). POP3 --- Post Office Protocol
version 3 --- downloads mail to the local machine; IMAP keeps mail synced on
the server (TRAP-17). Secure transmission uses SSL/TLS. Decimal 10 equals
binary 1010.

\section{Networking, security and cloud ladder}

BGP is an inter-domain routing protocol used between autonomous systems; OSPF
and RIP are intra-domain (TRAP-20). ICMP is a control and diagnostic protocol
that carries no application data and uses no ports (TRAP-19). HTTP/3 uses QUIC
over UDP, not TCP (TRAP-21). TCP provides ordering with no timing guarantee;
UDP is connectionless. A virus needs a host program, a worm self-propagates
across networks, and a Trojan poses as useful software. Encryption provides
confidentiality while a digital signature provides authenticity and integrity
(TRAP-15). In the IaaS responsibility model, the provider manages the
infrastructure (networking and storage hardware) and the customer manages the
OS, runtime, applications, and data (TRAP-22).

\section{IT Act and e-governance ladder}

The Information Technology Act, 2000 is India's primary cyber law. Section 43
covers unauthorised access and data damage (penalty and compensation), and
Section 66 covers computer-related offences. The e-governance service models
are G2C (Government to Citizen), G2B (Government to Business), G2G
(Government to Government), and G2E (Government to Employees). The anchor
programmes are Digital India, NeGP/e-Kranti, and the NeGD (National
e-Governance Division). Keep the Section 43/66 pair and the G2C/G2B pair
straight: 43 is unauthorised access, 66 is computer offences; G2C serves the
citizen, G2B serves business.

\begin{definitionbox}{G2x models}
G2C delivers services to citizens, G2B to businesses, G2G between government
departments, and G2E to government employees.
\end{definitionbox}

\section{Exam retrieval and common confusions}

Ten distinctions prevent most errors across the whole paper.
\begin{enumerate}
  \item RAM is volatile; ROM is non-volatile and read-only (TRAP-01, TRAP-02).
  \item Cache is faster than RAM and is not a backup; the CPU never touches
    secondary storage directly (TRAP-03).
  \item Plotter, drum printer, and sound card are output; scanner, light pen,
    and digitizer are input (TRAP-12).
  \item A device driver is software; BIOS boots the machine (TRAP-10,
    TRAP-24).
  \item Section Break changes headers; Page Break does not (TRAP-13).
  \item Formulas start with \texttt{=}; \texttt{\$A\$1} is absolute;
    Select query retrieves (TRAP-30, TRAP-31).
  \item F5 runs the show from the beginning; F7 checks spelling (TRAP-32).
  \item IPv6 is 128-bit; \texttt{@} separates user name and domain name
    (TRAP-18, TRAP-28).
  \item POP3 downloads locally; BCC hides recipients from each other
    (TRAP-16, TRAP-17).
  \item Signature is not encryption; BGP is inter-domain; HTTP/3 is QUIC over
    UDP; IaaS provider manages hardware only (TRAP-15, TRAP-20, TRAP-21,
    TRAP-22).
\end{enumerate}

\begin{example}
An item states three claims: cache is faster than RAM; POP3 downloads mail to
the local machine; IPv6 addresses are 256-bit. The first two trace to the
ladder and are true; the third inverts TRAP-18 (IPv6 is 128-bit) and is false.
Eliminate any option that asserts all three.
\end{example}

\subsection*{Exercises}
\begin{enumerate}[label=\arabic*.]
  \item State the memory speed order and the volatility of RAM and ROM.
  \item Classify plotter, scanner, sound card, and digitizer tablet as input
    or output, with one line of reasoning each.
  \item Distinguish Section Break from Page Break and F5 from F7.
  \item Explain the meaning of \texttt{=}, \texttt{\$A\$1}, and COUNTIFS in
    one line each.
  \item Contrast POP3 with IMAP and BCC with Cc.
  \item State the roles of BGP, ICMP, and QUIC in one line each.
  \item Expand G2C, G2B, G2G, and G2E, and name the Act and sections behind
    Sections 43 and 66.
\end{enumerate}

\subsection*{Solutions}
\begingroup\small
\begin{enumerate}[label=\arabic*.]
  \item Registers, cache (L1/L2/L3), RAM, secondary storage, tertiary/offline.
    RAM is volatile; ROM is non-volatile and read-only in normal operation.
  \item Plotter: output (draws); scanner: input (reads in); sound card: output
    (plays sound); digitizer tablet: input (captures strokes).
  \item Section Break restarts headers/footers; Page Break only starts a new
    page. F5 starts the slide show from the beginning; F7 runs spell check.
  \item \texttt{=} marks a formula; \texttt{\$A\$1} fixes both column and row;
    COUNTIFS counts rows meeting several criteria.
  \item POP3 downloads mail locally; IMAP syncs with the server. BCC hides
    recipients from each other; Cc stays visible to all recipients.
  \item BGP routes between autonomous systems; ICMP carries control and
    diagnostic messages, not data; QUIC over UDP carries HTTP/3.
  \item G2C: to Citizen; G2B: to Business; G2G: to Government; G2E: to
    Employees. Information Technology Act, 2000: Section 43, unauthorised
    access and data damage; Section 66, computer-related offences.
\end{enumerate}
\endgroup
\end{document}
